\documentclass[11pt]{article}
\usepackage{../common/ai204}
\renewcommand{\doctitle}{BTL Question Bank}
\begin{document}
\aicover{Question Bank}{30 questions per unit \textbullet\ 150 questions mapped to Bloom's Taxonomy (L1--L6) and Course Outcomes}{Remember \textbullet\ Understand \textbullet\ Apply \textbullet\ Analyse \textbullet\ Evaluate \textbullet\ Create}
\section*{How to use this Question Bank}
Every unit has exactly \textbf{30 questions}: five at each Bloom's Taxonomy Level (BTL). Marks guide: L1 = 2 marks, L2--L3 = 5 marks, L4--L6 = 10 marks. Use L1--L3 for CIA-I/quiz, L3--L5 for CIA-II and end-semester, L6 for assignments and open-book tasks.
\begin{center}\begin{tabular}{cllp{8cm}}
\rowcolor{navy}\textcolor{gold}{\textbf{Level}}&\textcolor{gold}{\textbf{Name}}&\textcolor{gold}{\textbf{Marks}}&\textcolor{gold}{\textbf{Typical verbs}}\\
\btl{L1}&Remember&2&define, list, state, name, recall\\
\rowcolor{ivory}\btl{L2}&Understand&5&explain, describe, discuss, summarise\\
\btl{L3}&Apply&5&apply, compute, solve, trace, formulate\\
\rowcolor{ivory}\btl{L4}&Analyse&10&analyse, compare, differentiate, examine\\
\btl{L5}&Evaluate&10&evaluate, justify, assess, critique, judge\\
\rowcolor{ivory}\btl{L6}&Create&10&design, develop, propose, construct, create\\
\end{tabular}\end{center}
\begin{center}\begin{tikzpicture}[x=1.7cm,y=0.18cm]
\draw[slate,->](0,0)--(6.6,0); \draw[slate,->](0,0)--(0,30) node[above,font=\small]{questions};
\fill[teal] (0.2,0) rectangle (0.8,25); \node[font=\small\bfseries] at (0.5,26.8){25}; \node[below,font=\small] at (0.5,0){L1};
\fill[navyL] (1.2,0) rectangle (1.8,25); \node[font=\small\bfseries] at (1.5,26.8){25}; \node[below,font=\small] at (1.5,0){L2};
\fill[leaf] (2.2,0) rectangle (2.8,25); \node[font=\small\bfseries] at (2.5,26.8){25}; \node[below,font=\small] at (2.5,0){L3};
\fill[gold!85!black] (3.2,0) rectangle (3.8,25); \node[font=\small\bfseries] at (3.5,26.8){25}; \node[below,font=\small] at (3.5,0){L4};
\fill[crimson] (4.2,0) rectangle (4.8,25); \node[font=\small\bfseries] at (4.5,26.8){25}; \node[below,font=\small] at (4.5,0){L5};
\fill[navy] (5.2,0) rectangle (5.8,25); \node[font=\small\bfseries] at (5.5,26.8){25}; \node[below,font=\small] at (5.5,0){L6};
\end{tikzpicture}\\{\small\color{slate}Course-wide distribution: 25 questions at each BTL level (150 total).}\end{center}
\newpage
\unitbanner{I}{Introduction to AI and Intelligent Agents}{30 Questions \textbullet\ CO204.1}
\begin{longtable}{>{\bfseries\color{navy}}p{0.8cm}p{11.6cm}ccc}
\rowcolor{navy}\textcolor{gold}{Q.}&\textcolor{gold}{Question}&\textcolor{gold}{BTL}&\textcolor{gold}{CO}&\textcolor{gold}{M}\\\endhead
1&Define Artificial Intelligence. State any two definitions from the four categories proposed by Russell and Norvig.&\btl{L1}&\co{CO204.1}&2\\
\rowcolor{ivory}2&What is the Turing Test? List the four capabilities a machine needs to pass it.&\btl{L1}&\co{CO204.1}&2\\
3&Define an agent, a percept sequence and an agent function.&\btl{L1}&\co{CO204.1}&2\\
\rowcolor{ivory}4&Expand PEAS and state what each component specifies.&\btl{L1}&\co{CO204.1}&2\\
5&List the five components of a well-defined search problem.&\btl{L1}&\co{CO204.1}&2\\
\arrayrulecolor{gold}\midrule\arrayrulecolor{black}
\rowcolor{ivory}6&Explain the four approaches to AI (thinking/acting, humanly/rationally) with one example each.&\btl{L2}&\co{CO204.1}&5\\
7&Describe the major milestones in the history and evolution of AI from 1943 to the present.&\btl{L2}&\co{CO204.1}&5\\
\rowcolor{ivory}8&Explain the concept of rationality. Why is a rational agent not the same as an omniscient agent?&\btl{L2}&\co{CO204.1}&5\\
9&Discuss the contributions of mathematics, economics, neuroscience and psychology to the foundations of AI.&\btl{L2}&\co{CO204.1}&5\\
\rowcolor{ivory}10&Explain the structure of a learning agent with a neat diagram, describing the role of the critic and problem generator.&\btl{L2}&\co{CO204.1}&5\\
\arrayrulecolor{gold}\midrule\arrayrulecolor{black}
11&Write the PEAS description for (a) an automated taxi driver and (b) an online English tutor.&\btl{L3}&\co{CO204.1}&5\\
\rowcolor{ivory}12&Classify the task environment of (a) crossword puzzle, (b) poker and (c) medical diagnosis along all seven properties.&\btl{L3}&\co{CO204.1}&5\\
13&Formulate the 8-puzzle problem as a state-space search problem. How many states are reachable?&\btl{L3}&\co{CO204.1}&5\\
\rowcolor{ivory}14&Formulate the Water-Jug problem (4 L and 3 L jugs, goal 2 L) and give one solution path.&\btl{L3}&\co{CO204.1}&5\\
15&Write the agent function in tabular form for the two-square vacuum-cleaner world.&\btl{L3}&\co{CO204.1}&5\\
\arrayrulecolor{gold}\midrule\arrayrulecolor{black}
\rowcolor{ivory}16&Compare simple-reflex, model-based, goal-based and utility-based agents in terms of internal state, knowledge required and flexibility.&\btl{L4}&\co{CO204.1}&10\\
17&Analyse why a simple-reflex agent may loop forever in a partially observable environment and how randomisation or a model helps.&\btl{L4}&\co{CO204.1}&10\\
\rowcolor{ivory}18&Distinguish between a state and a node in a search tree. Analyse how repeated states inflate the Water-Jug search tree.&\btl{L4}&\co{CO204.1}&10\\
19&Differentiate episodic vs sequential and static vs dynamic environments with examples from Indian application domains (UPI fraud detection, traffic signals).&\btl{L4}&\co{CO204.1}&10\\
\rowcolor{ivory}20&Analyse the Missionaries-and-Cannibals problem: identify states, legal actions and illegal states.&\btl{L4}&\co{CO204.1}&10\\
\arrayrulecolor{gold}\midrule\arrayrulecolor{black}
21&Evaluate the Turing Test as a criterion of machine intelligence in the light of Searle's Chinese Room argument and modern LLMs.&\btl{L5}&\co{CO204.1}&10\\
\rowcolor{ivory}22&Critically assess whether a self-driving car should be designed as a goal-based or a utility-based agent. Justify your choice.&\btl{L5}&\co{CO204.1}&10\\
23&Judge the suitability of a performance measure 'amount of dirt cleaned' for a vacuum agent. Propose and defend a better measure.&\btl{L5}&\co{CO204.1}&10\\
\rowcolor{ivory}24&Evaluate the claim that 'any AI problem can be solved if it is formulated as a state-space search'. Give supporting and counter examples.&\btl{L5}&\co{CO204.1}&10\\
25&Assess the impact of AI applications in Indian healthcare and agriculture, citing benefits and risks.&\btl{L5}&\co{CO204.1}&10\\
\arrayrulecolor{gold}\midrule\arrayrulecolor{black}
\rowcolor{ivory}26&Design a complete PEAS specification and agent architecture for a smart campus energy-management agent at Sanjivani University.&\btl{L6}&\co{CO204.1}&10\\
27&Propose a learning-agent design for a personalised tutoring system; identify the performance element, critic, learning element and problem generator.&\btl{L6}&\co{CO204.1}&10\\
\rowcolor{ivory}28&Formulate a real-world exam-timetabling problem as a well-defined search problem, specifying all five components.&\btl{L6}&\co{CO204.1}&10\\
29&Develop a taxonomy chart that maps ten everyday AI applications to environment properties and suitable agent types.&\btl{L6}&\co{CO204.1}&10\\
\rowcolor{ivory}30&Create a state-space representation for a delivery-drone routing problem in a city, including action costs and goal test.&\btl{L6}&\co{CO204.1}&10\\
\end{longtable}
\newpage
\unitbanner{II}{Problem Solving by Search}{30 Questions \textbullet\ CO204.2}
\begin{longtable}{>{\bfseries\color{navy}}p{0.8cm}p{11.6cm}ccc}
\rowcolor{navy}\textcolor{gold}{Q.}&\textcolor{gold}{Question}&\textcolor{gold}{BTL}&\textcolor{gold}{CO}&\textcolor{gold}{M}\\\endhead
1&Define branching factor, depth of solution and maximum depth of a search tree.&\btl{L1}&\co{CO204.2}&2\\
\rowcolor{ivory}2&State the data structure used for the frontier in BFS, DFS and UCS.&\btl{L1}&\co{CO204.2}&2\\
3&Define an admissible heuristic and a consistent heuristic.&\btl{L1}&\co{CO204.2}&2\\
\rowcolor{ivory}4&What is a constraint satisfaction problem? State its three components.&\btl{L1}&\co{CO204.2}&2\\
5&State the evaluation function used by Greedy Best-First Search and A*.&\btl{L1}&\co{CO204.2}&2\\
\arrayrulecolor{gold}\midrule\arrayrulecolor{black}
\rowcolor{ivory}6&Explain Iterative-Deepening Search and why its overhead is small compared to BFS.&\btl{L2}&\co{CO204.2}&5\\
7&Explain the problems of local maxima, ridges and plateaux in hill climbing and the remedies.&\btl{L2}&\co{CO204.2}&5\\
\rowcolor{ivory}8&Describe simulated annealing and the role of temperature in accepting worse moves.&\btl{L2}&\co{CO204.2}&5\\
9&Explain the working of a genetic algorithm with selection, crossover and mutation for the 8-queens problem.&\btl{L2}&\co{CO204.2}&5\\
\rowcolor{ivory}10&Explain bidirectional search and state when it cannot be applied.&\btl{L2}&\co{CO204.2}&5\\
\arrayrulecolor{gold}\midrule\arrayrulecolor{black}
11&Apply BFS and DFS to a given graph (A--G) and write the order of node expansion and the path found.&\btl{L3}&\co{CO204.2}&5\\
\rowcolor{ivory}12&Apply A* on the Romania map from Arad to Bucharest using straight-line distance and show f, g, h for each expansion.&\btl{L3}&\co{CO204.2}&5\\
13&Compute h1 (misplaced tiles) and h2 (Manhattan distance) for a given 8-puzzle start state.&\btl{L3}&\co{CO204.2}&5\\
\rowcolor{ivory}14&Apply Minimax to a given 3-ply game tree and find the value of the root and the best move.&\btl{L3}&\co{CO204.2}&5\\
15&Solve the Australia map-colouring CSP using backtracking with MRV and forward checking.&\btl{L3}&\co{CO204.2}&5\\
\arrayrulecolor{gold}\midrule\arrayrulecolor{black}
\rowcolor{ivory}16&Compare BFS, DFS, DLS, IDS, UCS and bidirectional search on completeness, optimality, time and space complexity.&\btl{L4}&\co{CO204.2}&10\\
17&Analyse why A* with a consistent heuristic never re-expands a node, while an inadmissible heuristic may return a sub-optimal path.&\btl{L4}&\co{CO204.2}&10\\
\rowcolor{ivory}18&Apply alpha--beta pruning to a given game tree; identify pruned branches and analyse how move ordering affects pruning.&\btl{L4}&\co{CO204.2}&10\\
19&Differentiate global search (A*) from local search (hill climbing) with respect to memory, optimality and application domains.&\btl{L4}&\co{CO204.2}&10\\
\rowcolor{ivory}20&Analyse AC-3 on a small CSP and show the domain reductions after each arc revision.&\btl{L4}&\co{CO204.2}&10\\
\arrayrulecolor{gold}\midrule\arrayrulecolor{black}
21&Evaluate which search strategy is best suited for a GPS navigation app with millions of nodes. Justify with complexity arguments.&\btl{L5}&\co{CO204.2}&10\\
\rowcolor{ivory}22&Judge whether the heuristic h(n) = 2 $\times$ Manhattan distance is acceptable for the 8-puzzle; discuss the speed--optimality trade-off.&\btl{L5}&\co{CO204.2}&10\\
23&Critically compare simulated annealing and genetic algorithms for solving the Travelling Salesperson Problem.&\btl{L5}&\co{CO204.2}&10\\
\rowcolor{ivory}24&Assess the limitations of Minimax for real games such as chess and Go and evaluate how MCTS overcomes them.&\btl{L5}&\co{CO204.2}&10\\
25&Evaluate the effectiveness of MRV, degree and LCV heuristics on Sudoku solving.&\btl{L5}&\co{CO204.2}&10\\
\arrayrulecolor{gold}\midrule\arrayrulecolor{black}
\rowcolor{ivory}26&Design an admissible heuristic for a robot path-planning problem on a grid with diagonal moves and prove its admissibility.&\btl{L6}&\co{CO204.2}&10\\
27&Develop a complete Tic-Tac-Toe agent using Minimax with alpha--beta pruning; write the pseudocode and evaluation function.&\btl{L6}&\co{CO204.2}&10\\
\rowcolor{ivory}28&Formulate the university exam-scheduling problem as a CSP and propose a solving strategy.&\btl{L6}&\co{CO204.2}&10\\
29&Design a genetic algorithm (encoding, fitness, operators) for a class-timetable optimisation problem.&\btl{L6}&\co{CO204.2}&10\\
\rowcolor{ivory}30&Create a comparative experiment plan to benchmark BFS, IDS and A* on randomly generated 8-puzzle instances.&\btl{L6}&\co{CO204.2}&10\\
\end{longtable}
\newpage
\unitbanner{III}{Knowledge Representation and Reasoning}{30 Questions \textbullet\ CO204.3}
\begin{longtable}{>{\bfseries\color{navy}}p{0.8cm}p{11.6cm}ccc}
\rowcolor{navy}\textcolor{gold}{Q.}&\textcolor{gold}{Question}&\textcolor{gold}{BTL}&\textcolor{gold}{CO}&\textcolor{gold}{M}\\\endhead
1&What is a knowledge base? State the two operations of a knowledge-based agent.&\btl{L1}&\co{CO204.3}&2\\
\rowcolor{ivory}2&Define entailment, validity and satisfiability.&\btl{L1}&\co{CO204.3}&2\\
3&List the connectives of propositional logic in order of precedence.&\btl{L1}&\co{CO204.3}&2\\
\rowcolor{ivory}4&Define unification and the most general unifier.&\btl{L1}&\co{CO204.3}&2\\
5&What is a frame? Name its main components.&\btl{L1}&\co{CO204.3}&2\\
\arrayrulecolor{gold}\midrule\arrayrulecolor{black}
\rowcolor{ivory}6&Explain the architecture of a knowledge-based agent with the KB-AGENT algorithm.&\btl{L2}&\co{CO204.3}&5\\
7&Explain the difference between propositional logic and first-order logic with examples.&\btl{L2}&\co{CO204.3}&5\\
\rowcolor{ivory}8&Explain universal and existential quantifiers and the duality between them.&\btl{L2}&\co{CO204.3}&5\\
9&Describe semantic networks with an example showing inheritance through is-a links.&\btl{L2}&\co{CO204.3}&5\\
\rowcolor{ivory}10&Explain what an ontology is and describe the role of RDF and OWL.&\btl{L2}&\co{CO204.3}&5\\
\arrayrulecolor{gold}\midrule\arrayrulecolor{black}
11&Construct the truth table to prove that (P $\Rightarrow$ Q) $\equiv$ ($\neg$P $\vee$ Q).&\btl{L3}&\co{CO204.3}&5\\
\rowcolor{ivory}12&Convert the sentence (A $\Leftrightarrow$ B) $\Rightarrow$ C to Conjunctive Normal Form, showing each step.&\btl{L3}&\co{CO204.3}&5\\
13&Translate into FOL: (a) All students who study pass. (b) Some birds cannot fly. (c) Every person has a mother.&\btl{L3}&\co{CO204.3}&5\\
\rowcolor{ivory}14&Find the MGU of Knows(John, x) and Knows(y, Mother(y)); and of P(x, f(y)) and P(a, f(g(z))).&\btl{L3}&\co{CO204.3}&5\\
15&Represent the Hotel-Room domain using frames with slots, default values and procedures.&\btl{L3}&\co{CO204.3}&5\\
\arrayrulecolor{gold}\midrule\arrayrulecolor{black}
\rowcolor{ivory}16&Apply forward chaining to the Colonel West crime knowledge base and show each iteration.&\btl{L4}&\co{CO204.3}&10\\
17&Apply backward chaining to prove Criminal(West) and draw the proof tree.&\btl{L4}&\co{CO204.3}&10\\
\rowcolor{ivory}18&Compare forward and backward chaining with respect to direction, efficiency and application domains.&\btl{L4}&\co{CO204.3}&10\\
19&Analyse the common errors in translating English to FOL involving $\forall$ with $\wedge$ and $\exists$ with $\Rightarrow$.&\btl{L4}&\co{CO204.3}&10\\
\rowcolor{ivory}20&Distinguish between semantic networks, frames and ontologies on expressiveness and reasoning support.&\btl{L4}&\co{CO204.3}&10\\
\arrayrulecolor{gold}\midrule\arrayrulecolor{black}
21&Prove using resolution: Given $\forall$x Man(x) $\Rightarrow$ Mortal(x) and Man(Socrates), prove Mortal(Socrates); justify each step.&\btl{L5}&\co{CO204.3}&10\\
\rowcolor{ivory}22&Evaluate the soundness and completeness of resolution as an inference procedure.&\btl{L5}&\co{CO204.3}&10\\
23&Assess the limitations of propositional logic for representing the Wumpus World and justify the need for FOL.&\btl{L5}&\co{CO204.3}&10\\
\rowcolor{ivory}24&Critically evaluate rule-based (logic) knowledge representation against neural representations used in LLMs.&\btl{L5}&\co{CO204.3}&10\\
25&Judge whether the sentence $\forall$x $\exists$y Loves(x,y) entails $\exists$y $\forall$x Loves(x,y). Justify with a model.&\btl{L5}&\co{CO204.3}&10\\
\arrayrulecolor{gold}\midrule\arrayrulecolor{black}
\rowcolor{ivory}26&Design a family-relationship knowledge base in FOL (parent, grandparent, sibling, cousin) and pose three queries.&\btl{L6}&\co{CO204.3}&10\\
27&Develop a Prolog program for a university course-prerequisite knowledge base and trace a query.&\btl{L6}&\co{CO204.3}&10\\
\rowcolor{ivory}28&Create an ontology (classes, properties, individuals) for the AI \& DS department using RDF triples.&\btl{L6}&\co{CO204.3}&10\\
29&Construct a complete resolution refutation proof for a given knowledge base about 'Marcus' (Rich \& Knight).&\btl{L6}&\co{CO204.3}&10\\
\rowcolor{ivory}30&Design a semantic network for an Indian railway reservation system showing inheritance and attributes.&\btl{L6}&\co{CO204.3}&10\\
\end{longtable}
\newpage
\unitbanner{IV}{Planning and Reasoning under Uncertainty}{30 Questions \textbullet\ CO204.4}
\begin{longtable}{>{\bfseries\color{navy}}p{0.8cm}p{11.6cm}ccc}
\rowcolor{navy}\textcolor{gold}{Q.}&\textcolor{gold}{Question}&\textcolor{gold}{BTL}&\textcolor{gold}{CO}&\textcolor{gold}{M}\\\endhead
1&State the components of a planning system.&\btl{L1}&\co{CO204.4}&2\\
\rowcolor{ivory}2&What does STRIPS stand for? List the parts of a STRIPS action schema.&\btl{L1}&\co{CO204.4}&2\\
3&State Bayes' theorem and name each term.&\btl{L1}&\co{CO204.4}&2\\
\rowcolor{ivory}4&Define a Bayesian network and a conditional probability table.&\btl{L1}&\co{CO204.4}&2\\
5&Define a fuzzy set and a membership function.&\btl{L1}&\co{CO204.4}&2\\
\arrayrulecolor{gold}\midrule\arrayrulecolor{black}
\rowcolor{ivory}6&Explain forward (progression) and backward (regression) state-space planning.&\btl{L2}&\co{CO204.4}&5\\
7&Explain partial-order planning with causal links, ordering constraints and threats.&\btl{L2}&\co{CO204.4}&5\\
\rowcolor{ivory}8&Describe the sources of uncertainty faced by an intelligent agent.&\btl{L2}&\co{CO204.4}&5\\
9&Explain conditional independence and how it makes Bayesian networks compact.&\btl{L2}&\co{CO204.4}&5\\
\rowcolor{ivory}10&Explain the steps of a Mamdani fuzzy inference system with a block diagram.&\btl{L2}&\co{CO204.4}&5\\
\arrayrulecolor{gold}\midrule\arrayrulecolor{black}
11&Write STRIPS operators for PickUp, PutDown, Stack and Unstack in the Blocks World.&\btl{L3}&\co{CO204.4}&5\\
\rowcolor{ivory}12&A disease has prevalence 1\%; a test has 90\% sensitivity and 5\% false-positive rate. Compute P(Disease | +).&\btl{L3}&\co{CO204.4}&5\\
13&Using the burglary network, compute P(j, m, a, $\neg$b, $\neg$e).&\btl{L3}&\co{CO204.4}&5\\
\rowcolor{ivory}14&Draw triangular membership functions for Cold, Warm and Hot temperature and find the membership of 18 $^\circ$C.&\btl{L3}&\co{CO204.4}&5\\
15&Compute P(cavity | toothache) from the given full joint distribution table.&\btl{L3}&\co{CO204.4}&5\\
\arrayrulecolor{gold}\midrule\arrayrulecolor{black}
\rowcolor{ivory}16&Generate a plan for a given Blocks World problem using goal-stack planning and analyse the Sussman anomaly.&\btl{L4}&\co{CO204.4}&10\\
17&Compare total-order and partial-order planning with an example (putting on shoes and socks).&\btl{L4}&\co{CO204.4}&10\\
\rowcolor{ivory}18&Compare probability theory and fuzzy logic for handling uncertainty and vagueness, with examples.&\btl{L4}&\co{CO204.4}&10\\
19&Analyse inference by enumeration for P(Burglary | JohnCalls, MaryCalls) and comment on its complexity.&\btl{L4}&\co{CO204.4}&10\\
\rowcolor{ivory}20&Apply the fuzzy operators union, intersection, complement and the hedge 'very' on two given fuzzy sets.&\btl{L4}&\co{CO204.4}&10\\
\arrayrulecolor{gold}\midrule\arrayrulecolor{black}
21&Evaluate the base-rate fallacy using the medical test example and its implications for AI diagnosis tools.&\btl{L5}&\co{CO204.4}&10\\
\rowcolor{ivory}22&Assess the assumptions of classical planning and judge whether they hold for a Mars rover.&\btl{L5}&\co{CO204.4}&10\\
23&Critically evaluate Na\"{i}ve Bayes for spam filtering: when does the independence assumption fail?&\btl{L5}&\co{CO204.4}&10\\
\rowcolor{ivory}24&Judge whether exact or approximate inference should be used in a 100-node medical Bayesian network.&\btl{L5}&\co{CO204.4}&10\\
25&Evaluate the effect of different defuzzification methods (centroid, mean-of-max) on a fan-speed controller.&\btl{L5}&\co{CO204.4}&10\\
\arrayrulecolor{gold}\midrule\arrayrulecolor{black}
\rowcolor{ivory}26&Design STRIPS/PDDL operators and a plan for a robot that delivers parcels between rooms.&\btl{L6}&\co{CO204.4}&10\\
27&Construct a Bayesian network for diagnosing a laptop that does not boot; give the structure and CPTs.&\btl{L6}&\co{CO204.4}&10\\
\rowcolor{ivory}28&Design a complete fuzzy controller (inputs, rules, outputs) for an automatic washing machine.&\btl{L6}&\co{CO204.4}&10\\
29&Propose a partial-order plan for preparing tea and identify all causal links and threats.&\btl{L6}&\co{CO204.4}&10\\
\rowcolor{ivory}30&Develop a probabilistic model for a crop-disease advisory system for farmers in Maharashtra.&\btl{L6}&\co{CO204.4}&10\\
\end{longtable}
\newpage
\unitbanner{V}{Learning, Expert Systems and Applications of AI}{30 Questions \textbullet\ CO204.5}
\begin{longtable}{>{\bfseries\color{navy}}p{0.8cm}p{11.6cm}ccc}
\rowcolor{navy}\textcolor{gold}{Q.}&\textcolor{gold}{Question}&\textcolor{gold}{BTL}&\textcolor{gold}{CO}&\textcolor{gold}{M}\\\endhead
1&Define machine learning as given by Tom Mitchell.&\btl{L1}&\co{CO204.5}&2\\
\rowcolor{ivory}2&List the three forms of learning with one example each.&\btl{L1}&\co{CO204.5}&2\\
3&Define entropy and information gain.&\btl{L1}&\co{CO204.5}&2\\
\rowcolor{ivory}4&What is a perceptron? Write its output equation.&\btl{L1}&\co{CO204.5}&2\\
5&List the components of an expert system.&\btl{L1}&\co{CO204.5}&2\\
\arrayrulecolor{gold}\midrule\arrayrulecolor{black}
\rowcolor{ivory}6&Explain supervised, unsupervised and reinforcement learning with suitable examples.&\btl{L2}&\co{CO204.5}&5\\
7&Explain how a decision tree is learned using the ID3 algorithm.&\btl{L2}&\co{CO204.5}&5\\
\rowcolor{ivory}8&Explain the perceptron learning rule and the perceptron convergence theorem.&\btl{L2}&\co{CO204.5}&5\\
9&Explain the architecture of an expert system with a neat diagram.&\btl{L2}&\co{CO204.5}&5\\
\rowcolor{ivory}10&Explain the knowledge acquisition process and why it is called a bottleneck.&\btl{L2}&\co{CO204.5}&5\\
\arrayrulecolor{gold}\midrule\arrayrulecolor{black}
11&Compute the entropy of a dataset with 9 positive and 5 negative examples.&\btl{L3}&\co{CO204.5}&5\\
\rowcolor{ivory}12&Compute the information gain of attribute Outlook for the Play Tennis dataset.&\btl{L3}&\co{CO204.5}&5\\
13&Train a perceptron for the AND function ($\eta$ = 1, initial weights 0) for two epochs; show all updates.&\btl{L3}&\co{CO204.5}&5\\
\rowcolor{ivory}14&Apply forward chaining on a given medical rule base with certainty factors and compute the final CF.&\btl{L3}&\co{CO204.5}&5\\
15&Illustrate the stages of an NLP pipeline for the sentence 'Sanjivani University is in Kopargaon'.&\btl{L3}&\co{CO204.5}&5\\
\arrayrulecolor{gold}\midrule\arrayrulecolor{black}
\rowcolor{ivory}16&Analyse why a single perceptron cannot learn XOR and how a multilayer network solves it.&\btl{L4}&\co{CO204.5}&10\\
17&Compare decision trees and neural networks on interpretability, data requirement and accuracy.&\btl{L4}&\co{CO204.5}&10\\
\rowcolor{ivory}18&Compare forward-chaining and backward-chaining expert systems with examples (XCON, MYCIN).&\btl{L4}&\co{CO204.5}&10\\
19&Analyse overfitting in decision trees and compare pre-pruning and post-pruning.&\btl{L4}&\co{CO204.5}&10\\
\rowcolor{ivory}20&Differentiate the AI techniques used in NLP, computer vision and robotics.&\btl{L4}&\co{CO204.5}&10\\
\arrayrulecolor{gold}\midrule\arrayrulecolor{black}
21&Evaluate the ethical issues of bias and fairness in AI-based hiring or loan approval systems.&\btl{L5}&\co{CO204.5}&10\\
\rowcolor{ivory}22&Assess the advantages and limitations of expert systems compared to machine-learning systems.&\btl{L5}&\co{CO204.5}&10\\
23&Critically evaluate the use of facial recognition by public authorities in India with reference to privacy laws.&\btl{L5}&\co{CO204.5}&10\\
\rowcolor{ivory}24&Judge the need for explainability (XAI) in medical AI and evaluate LIME/SHAP-style explanations.&\btl{L5}&\co{CO204.5}&10\\
25&Evaluate the EU AI Act's risk-based approach and its relevance to Indian AI policy.&\btl{L5}&\co{CO204.5}&10\\
\arrayrulecolor{gold}\midrule\arrayrulecolor{black}
\rowcolor{ivory}26&Design a rule-based expert system (at least 10 rules) for crop-disease diagnosis and show a sample consultation.&\btl{L6}&\co{CO204.5}&10\\
27&Build a decision tree by hand for a given student-placement dataset and extract IF--THEN rules.&\btl{L6}&\co{CO204.5}&10\\
\rowcolor{ivory}28&Propose a responsible-AI checklist for a student-performance prediction system at a university.&\btl{L6}&\co{CO204.5}&10\\
29&Design an AI-based chatbot for college admissions: architecture, NLP components and evaluation plan.&\btl{L6}&\co{CO204.5}&10\\
\rowcolor{ivory}30&Develop a reinforcement-learning formulation (states, actions, rewards) for a smart traffic-signal controller.&\btl{L6}&\co{CO204.5}&10\\
\end{longtable}
\newpage
\end{document}